\section{Lean formalization}\label{sec:lean-formalization}

\hy{Make sure that we use the most recent version of Lean 4 and Mathlib. The old versions
have a bug. Also, the repository should be different from our work repository. Some statistics
about formalisation would be useful. Axioms assumed need to be explained clearly. Also,
when preparing the public repository, make sure that we clean up the Lean code by refactoring
and removing dead code etc, make sure that the comments inside the code are consistent with
the paper, and also that the documentations, i.e., md files, are clear and up to date. We
can do most of these using Claude or Codex. We also have to double-check the statements
of theorems and definitions in the Lean code. Finally, there is a good chance that
Lean formalization deviates from the paper in some places. We need to try to minimize this,
but if it happens, we should clearly explain the differences in the documentation.}
A substantial part of the paper has been formalized in Lean 4 using Mathlib;
the code is available in the
\href{https://github.com/Lim-Sangho/upper-tail-optimizers-lean}{online
repository}.  It contains no unproved \texttt{sorry} placeholders.  Its
purpose is to machine-check the deterministic graphon arguments underlying our
results.  It does not formalize the Chatterjee--Varadhan large deviation
principle, and hence does not cover the probabilistic consequences stated in
the introduction.

\paragraph{Nonexceptional theory.}
The main results away from the exceptional density are formalized, ending in
\Cref{thm:nonexceptional-optimizers}.  This includes the construction and analysis
of the Lubetzky--Zhao boundary in
\Cref{sec:lz-boundary}, the reduction to a one-dimensional problem in
\Cref{sec:local-reduction}, the positivity of $A_H$ in
\Cref{sec:nondegeneracy}, the uniqueness and analyticity results and the
edge-density and rate-function expansions proved in
\Cref{sec:local-optimizer-structure}, and the parameter expansions derived in
Appendix~\ref{app:bipodal-parameter-expansions}.
Appendix~\ref{app:krrs-analytic-extension}
is also formalized, so the two-sided extension of the KRR--S parametrization,
\Cref{thm:krrs-analytic-extension}, is proved within the formal development rather than assumed.
Two qualifications.  This part of the formalization treats the $d$-regular case
needed for our main results, rather than the more general $d$-starlike
formulation of \Cref{thm:krrs-analytic-extension}.  And the replica-symmetric half of
\Cref{thm:local-optimizer-structure} is formalized on the window $\pc(r)\le p<p_*$ rather than
on the full range $\pc(r)\le p<r$ for which we state it; this has no effect on
\Cref{thm:local-optimizer-structure} itself, whose neighborhood already forces $p<p_*$.

\paragraph{Exceptional endpoint.}
\Cref{sec:singular-endpoint} is formalized, including its terminal
theorem.  The analytic construction of the rank-one family converging to the
constant endpoint, its parameter expansions and its strict improvement over
the constant graphon, the localization and rank-one reduction, the rank-one
factor comparison, and the arbitrary-graphon comparison are all
machine-checked, and \Cref{thm:endpoint-optimality} is proved in both of its
clauses, global optimality and uniqueness up to measure-preserving relabelling.
Both conclusions include the construction of this family rather than
assuming its prior existence.  \Cref{thm:endpoint-optimizers} is
assembled from these results as a single formal statement.

\paragraph{External results.}
The only mathematical results assumed as axioms are the standard graphon
compactness, continuity, and lower-semicontinuity properties used in the paper;
the generalized H\"older inequality \eqref{eq:generalized-holder}; the
Lubetzky--Zhao criterion \Cref{thm:lz-criterion}; and two published
entropy-maximization results
of Kenyon, Radin, Ren, and Sadun
\cite[Theorems~1.1 and~3.3]{kenyon2014entropy}.  Each is assumed with the
standing hypotheses of its source, and all other statements described
above are derived from these results inside Lean.  One reading of the quoted
statements is worth recording: as in \Cref{sec:preliminaries} and
Appendix~\ref{app:krrs-analytic-extension}, we take the positive-surplus region of
\cite[Theorem~1.1]{kenyon2014entropy} to be open, since the analyticity of the
parameter maps on it presupposes as much, and that openness is what turns the
per-$\eps$ threshold $\tau_0(\eps)$ into the surplus window
\eqref{eq:krrs-uniform-domain} valid on a neighbourhood of $\eps_0$.  The
repository also contains a separate development, of the $p\downarrow0$ boundary
for triangles, which carries an assumption of its own; nothing described here
depends on it.